\section{Related Work}\label{sec:related}

The Peira framework we adopted in this study is part of a wider discourse in the research community on conceptualizing and evaluating model and modeling language qualities. While early frameworks for conceptually organizing the dimensions of quality, such as by Lindland et al. \cite{LiSS.94} had a strong focus on the quality of models, later proposals also explored appropriateness dimensions of the languages/grammars used to develop the models \cite{Krog.12,NePG.12,Mood.09,WaWe.93,Guiz.05}. The role of empirical studies, and particularly experiments, in assessing such qualities has also been recognized as evident by both the numerous studies published over the past decades (e.g., \cite{CGHM.13,LiHJ.13,HBKPRS.13,AlZL.17,LiRZ.17,LiTa.19} for models akin to iStar/iStar-DT alone) and the efforts to develop guidelines for systematically conducting such empirical evaluations and avoiding pitfalls \cite{BuWW.09,PaCo.05,KGPP.02,Mood.05,HoFL.12}. Peira is the outcome of a longer-term effort \cite{LiJa.20,LiMK.21b} to offer specific and readily executable data collection and analysis procedures and artifacts, based on metrics that are direct operationalizations of established theories of language quality.  

Since the advent of LLMs, research has emerged investigating their role in assisting with conceptual modeling tasks, primarily automated model generation. Studies have been conducted for a variety of languages including UML (e.g., \cite{CTBV.23,DGCT+24,FeAA.24,XiSB.25}), BPMN (e.g., \cite{HoMR.26,LPRM.25,NAAD+25,WWLL+24}) and goal models \cite{CCHY+23}. More relevant to our work is the investigation of LLMs in assisting with model quality. Gutschmidt and Nast \cite{GuNa.25}, for instance, use an LLM to evaluate aspects of model quality, referring to the SEQUAL quality framework \cite{Krog.12}, while Ayad and AlSayoud \cite{AyAl.24} explore the use of GPT-4 for eliciting business process model improvements. We could not find work comparing human and LLM performance in evaluations of language vocabularies using Peira or otherwise.
